\section{Design Principles}\label{sec:principles}

Eight rules shape the design; the rest of the report shows where each applies.

\emph{Delete only what has a checked second copy.} The original is deleted only after GitHub's
branch has been read back and shows every chunk under the id computed locally from its
ciphertext, a sample has been downloaded and compared, and the file is still the one that was
read. Remote chunks are deleted only after the restored file has been written, checked, flushed
to disk and given its final name.

\emph{Record the intent before acting.} Every step that changes something outside the data folder
is preceded by a durable record of what is about to happen. The record names everything the step
could create (chunk paths are chosen at random and recorded before the first upload), so nothing
can appear on GitHub that GhostVault could not find and remove.

\emph{A reply is not proof.} A reply can be lost, can come from a lagging replica, or can report
success for a change that did not take effect. Wherever an outcome matters, GhostVault reads the
branch again and acts on what it shows.

\emph{GitHub stores; it is not trusted with secrets or integrity.} Only ciphertext under keys
GitHub never sees is uploaded, and every byte that comes back is checked against its content id,
its SHA-256 and its authentication tag before it is used.

\emph{One random key per file.} File keys are random, never derived from a password, so a weak
passphrase can weaken only the recovery bundle it protects, and a key that leaks exposes one file.

\emph{Validate before acting.} A record is validated every time it is read. Every value used to
build a request, choose a path or delete something must have the only form GhostVault itself
creates, so a damaged or edited record cannot direct a deletion at anything that is not
GhostVault's own.

\emph{When in doubt, keep the data.} An unreadable record is reported and never overwritten; a
damaged settings file is moved aside, never deleted; an original that changed is kept; a backup
that fails never fails the operation that triggered it.

\emph{Say what is true.} Status shown to the user reflects what was verified. A ghost counts as
protected by the repository backup only once that version has been read back from GitHub, and a
file whose original could not be deleted is reported as kept, with the reason.
